% ROZDZIAL 30: miary wektorowe i wlasnosc Radona-Nikodyma
\chapter[Miary wektorowe i własność Radona--Nikodyma]{Miary wektorowe, całka Bochnera i własność Radona--Nikodyma}
\label{ch:vector-measures-30}

\section{Od miar skalarnych do miar w przestrzeni Banacha}
\label{sec:vector-intro-30}
W Rozdziale~\ref{ch:geometric-measure-29} pojawiła się
pochodna $Du$ o wartościach w $\mathbb R^n$,
a w Rozdziale~\ref{ch:area-coarea-28} różniczka
odwzorowania między przestrzeniami skończenie wymiarowymi.
W pierwszym przypadku dla $u\in BV(\Omega)$ mamy
skończoną wariację miary $Du$, natomiast dla
$u\in BV_{\mathrm{loc}}(\Omega)$ tylko wariację
lokalnie skończoną. Definicja miary wektorowej
poniżej wymaga, aby wartość na każdym zbiorze
mierzalnym należała do $X$; nie należy bez dodatkowych
założeń podstawiać do niej lokalnie skończonej
miary $Du$ na całym $\Omega$. Teraz pozwalamy, aby miara przyjmowała
wartości w \emph{dowolnej} rzeczywistej lub zespolonej przestrzeni
Banacha $X$. Szczególnie ważne okaże się pytanie, kiedy istnieje
funkcja $g:\Omega\to X$ taka, że $\nu(A)=\int_Ag\,\mathrm d\mu$.
W przypadku skalarnym odpowiedź daje twierdzenie Radona--Nikodyma;
w przypadku Banachowskim zależy ona od geometrii $X$.

\begin{definicja}[Miara wektorowa]
\label{def:vector-measure-30}
Niech $(\Omega,\mathcal A)$ będzie przestrzenią mierzalną,
a $X$ przestrzenią Banacha. \emph{Miarą wektorową o wartościach w $X$}
nazywamy odwzorowanie $\nu:\mathcal A\to X$, dla którego
$\nu(\varnothing)=0$ i dla każdej rodziny parami rozłącznych
$A_j\in\mathcal A$ zachodzi
\[
 \nu\left(\bigcup_{j=1}^{\infty}A_j\right)
 =\lim_{N\to\infty}\sum_{j=1}^{N}\nu(A_j)
 \quad\text{w normie }X.
\]
Przez $\|\nu(A)\|$ oznaczamy normę \emph{wektora} $\nu(A)$.
Nie należy jej mylić z wariacją miary na zbiorze $A$.
\end{definicja}

\begin{przyklad}[Dwie podstawowe konstrukcje]
\label{ex:basic-vector-measures-30}
Jeżeli $\mu$ jest dodatnią miarą skończoną, a $x_0\in X$,
to $\nu(A)=\mu(A)x_0$ jest miarą wektorową.
Jeśli $(\Omega,\mathcal A,\mu)$ jest przestrzenią o mierze
skończonej, a $g:\Omega\to X$ jest całkowalna w sensie
Bochnera (definicja poniżej), to
$\nu_g(A)=\int_Ag\,\mathrm d\mu$ jest miarą wektorową.
Zbieżność w normie sum po rozłącznych zbiorach
wynika z oszacowania
$\|\int_Bg\,\mathrm d\mu\|\le\int_B\|g\|\,\mathrm d\mu$
i bezwzględnej ciągłości całki skalarnej.
\end{przyklad}

\section{Wariacja całkowita i półwariacja}
\label{sec:vector-variation-30}
\begin{definicja}[Wariacja i półwariacja]
\label{def:vector-variation-30}
Niech $\nu:\mathcal A\to X$ będzie miarą wektorową.
Dla $A\in\mathcal A$ definiujemy \emph{wariację całkowitą}
\[
 |\nu|(A)=\sup\left\{\sum_{j=1}^{N}\|\nu(A_j)\|:
 A=\dot\bigcup_{j=1}^{N}A_j,\quad A_j\in\mathcal A,\ N\ge1\right\}.
\]
Jeżeli $|\nu|(\Omega)<\infty$, mówimy, że $\nu$ ma
\emph{skończoną wariację}. Dla $x^*\in X^*$ odwzorowanie
$A\mapsto x^*(\nu(A))$ jest miarą skalarną.
\emph{Półwariacją} nazywamy
$\|\nu\|_{\mathrm{sv}}(A)
=\sup_{\|x^*\|\le1}|x^*\nu|(A)$.
\end{definicja}

\begin{twierdzenie}[Wariacja jest miarą dodatnią]
\label{thm:vector-variation-measure-30}
Dla każdej miary wektorowej $\nu:\mathcal A\to X$
funkcja $|\nu|:\mathcal A\to[0,+\infty]$ jest miarą.
Ponadto dla każdego $A\in\mathcal A$ zachodzi
$\|\nu(A)\|\le\|\nu\|_{\mathrm{sv}}(A)\le|\nu|(A)$.
\end{twierdzenie}
\begin{proof}
Monotoniczność wariacji wynika z dopisania zbioru
$B\setminus A$ do każdego podziału $A\subseteq B$.
Niech $A,B$ będą rozłączne. Łącząc podziały $A$ i $B$,
otrzymujemy $|\nu|(A\cup B)\ge|\nu|(A)+|\nu|(B)$.
W drugą stronę dowolny skończony podział
$A\cup B=\dot\bigcup_i C_i$ rozcinamy przez $A$ i $B$.
Trójkąt daje
$\sum_i\|\nu(C_i)\|
\le\sum_i\|\nu(C_i\cap A)\|+\sum_i\|\nu(C_i\cap B)\|
\le|\nu|(A)+|\nu|(B)$.
Zatem wariacja jest skończenie addytywna, także wtedy,
gdy przyjmuje wartość nieskończoną.

Jeśli $A=\dot\bigcup_{j=1}^{\infty}A_j$,
to z monotoniczności i skończonej addytywności
$|\nu|(A)\ge\sum_j|\nu|(A_j)$.
Dla nierówności odwrotnej wybierzmy dowolny
skończony podział $A=\dot\bigcup_i C_i$.
Z przeliczalnej addytywności $\nu$ w normie wynika
\[
 \|\nu(C_i)\|\le\sum_j\|\nu(C_i\cap A_j)\|.
\]
Sumujemy po skończenie wielu $i$ i zamieniamy kolejność
sum nieujemnych, uzyskując
$\sum_i\|\nu(C_i)\|\le\sum_j|\nu|(A_j)$.
Supremum po podziałach daje nierówność odwrotną.
Dla dowolnego $x^*$ z $\|x^*\|\le1$ mamy
$|x^*\nu|(A)\le|\nu|(A)$. Supremum po $x^*$
daje drugą nierówność z tezy, zaś pierwsza wynika
z twierdzenia Hahna--Banacha:
$\|\nu(A)\|=\sup_{\|x^*\|\le1}|x^*(\nu(A))|$.
\end{proof}

\begin{przyklad}[Skończona półwariacja nie implikuje skończonej wariacji]
\label{ex:semivariation-not-variation-30}
Na $\Omega=\mathbb N$ z $\mathcal A=2^{\mathbb N}$
niech $X=\ell^2$ i $\nu(A)=\sum_{n\in A}e_n/n$,
gdzie $(e_n)$ jest bazą kanoniczną. Szereg po dowolnym
zbiorze zbiega bezwarunkowo w $\ell^2$, więc $\nu$
jest miarą wektorową. Dla podziału na singletony
$|\nu|(\mathbb N)\ge\sum_{n=1}^{N}1/n$; zatem
$|\nu|(\mathbb N)=+\infty$.
Natomiast dla $x^*=(a_n)\in\ell^2$ z $\|x^*\|\le1$
\[
 |x^*\nu|(\mathbb N)=\sum_{n=1}^{\infty}\frac{|a_n|}{n}
 \le\left(\sum_n|a_n|^2\right)^{1/2}
 \left(\sum_n\frac1{n^2}\right)^{1/2}.
\]
Równość w kresie górnym wynika z nierówności
Cauchy'ego--Schwarza i właściwego doboru $a_n$.
Półwariacja jest skończona, mimo że wariacja nie jest.
\end{przyklad}

\begin{lemat}[Bezwzględna ciągłość wariacji]
\label{lem:vector-variation-ac-30}
Jeśli $\nu$ ma skończoną wariację i $\nu\ll\mu$
dla dodatniej miary $\mu$, to $|\nu|\ll\mu$.
\end{lemat}
\begin{proof}
Jeśli $\mu(A)=0$, to $\nu(B)=0$ dla każdego
mierzalnego $B\subseteq A$. Wszystkie sumy definiujące
$|\nu|(A)$ są więc zerowe.
\end{proof}

\section{Całka Bochnera i miary z gęstością}
\label{sec:bochner-30}
\begin{definicja}[Silna mierzalność i całka Bochnera]
\label{def:bochner-30}
Niech $(\Omega,\mathcal A,\mu)$ będzie przestrzenią
z dodatnią miarą, a $X$ przestrzenią Banacha.
Funkcja $g:\Omega\to X$ jest \emph{silnie mierzalna},
jeżeli jest granicą punktową prawie wszędzie ciągu
funkcji prostych o wartościach w $X$.
Jest \emph{całkowalna w sensie Bochnera}, gdy dodatkowo
$\int_\Omega\|g(\omega)\|\,\mathrm d\mu(\omega)<\infty$.
Dla funkcji prostej $s=\sum_{j=1}^{N}x_j\chi_{A_j}$
o rozłącznych zbiorach $A_j$ o skończonej mierze
kładziemy $\int_\Omega s\,\mathrm d\mu=\sum_j\mu(A_j)x_j$.
Dla ogólnej $g$ całkę określa granica tych sum
w normie $X$ dla aproksymacji zbieżnej w
$L^1(\mu;X)$. Przestrzeń tych klas funkcji
oznaczamy $L^1(\mu;X)$.
\end{definicja}

\begin{twierdzenie}[Konstrukcja i oszacowanie całki Bochnera]
\label{thm:bochner-basic-30}
Przestrzeń $L^1(\mu;X)$ jest zupełna, a wzór
$g\mapsto\int_\Omega g\,\mathrm d\mu$ wyznacza
ciągły operator liniowy $L^1(\mu;X)\to X$
o normie co najwyżej $1$. Dla każdego $g\in L^1(\mu;X)$
zachodzi
\[
 \left\|\int_\Omega g\,\mathrm d\mu\right\|
 \le\int_\Omega\|g\|\,\mathrm d\mu.
\]
\end{twierdzenie}
\begin{proof}
Dla funkcji prostych nierówność jest zwykłą nierównością
trójkąta. Niech $g$ będzie silnie mierzalna i
$\int_\Omega\|g\|\,\mathrm d\mu<\infty$.
Z definicji silnej mierzalności wybierzmy proste
$s_n\to g$ prawie wszędzie. Możemy dobrać je tak, by
$\|s_n(\omega)\|\le2\|g(\omega)\|$ prawie wszędzie:
w punkcie, w którym pierwotna prosta jest zbyt duża,
zastępujemy jej wartość przez $0$. Otrzymany zbiór
wycięcia jest mierzalny, a po podziale na skończenie
wiele wartości prostej nadal otrzymujemy funkcję prostą.
Twierdzenie Lebesgue'a o zbieżności zdominowanej
zastosowane do norm daje $\|s_n-g\|_{L^1}\to0$.
Wybór przybliżeń z nośnikami o skończonej mierze
wymaga jeszcze uzasadnienia, gdy $\mu(\Omega)=\infty$.
Dla $k\ge1$ zbiór $\{\|g\|>1/k\}$ ma skończoną
miarę z nierówności Markowa. Ich suma pokrywa
$\{\|g\|>0\}$, a z dominowanej zbieżności wynika,
że $g\chi_{\{\|g\|>1/k\}}\to g$ w $L^1$.
Na każdym z tych zbiorów o skończonej mierze
powyższe funkcje proste można obcinać bez zmiany
argumentu dominacji. Zatem funkcje proste
o całkowalnych normach są gęste także bez
założenia $\sigma$-skończoności całej miary $\mu$.

Dla ciągu prostych Cauchy'ego w $L^1$ wybieramy
podciąg $s_{n_k}$ taki, że
$\sum_k\|s_{n_{k+1}}-s_{n_k}\|_{L^1}<\infty$.
Twierdzenie Beppo Leviego daje
$\sum_k\|s_{n_{k+1}}(\omega)-s_{n_k}(\omega)\|<\infty$
prawie wszędzie. Zupełność $X$ zapewnia punktową
zbieżność do silnie mierzalnej $g$; nierówność Fatou
daje $s_{n_k}\to g$ w $L^1$, a w konsekwencji zbiega
cały ciąg. To dowodzi zupełności. Nierówność dla całki
prostej i gęstość jednoznacznie przedłużają całkę
i oszacowanie na całe $L^1(\mu;X)$.
\end{proof}

\begin{twierdzenie}[Wariacja miary o gęstości Bochnera]
\label{thm:variation-density-30}
Niech $g\in L^1(\mu;X)$, a
$\nu_g(A)=\int_Ag\,\mathrm d\mu$ dla $A\in\mathcal A$.
Wtedy $\nu_g$ jest miarą wektorową o skończonej wariacji
oraz
\[
 |\nu_g|(A)=\int_A\|g(\omega)\|\,\mathrm d\mu(\omega)
 \quad(A\in\mathcal A).
\]
\end{twierdzenie}
\begin{proof}
Przeliczalna addytywność wynika z
$\|\nu_g(\bigcup_{j>N}A_j)\|
\le\int_{\bigcup_{j>N}A_j}\|g\|\,\mathrm d\mu\to0$.
Dla każdego podziału $A=\dot\bigcup_jA_j$ mamy
$\sum_j\|\nu_g(A_j)\|\le\int_A\|g\|\,\mathrm d\mu$.
W drugą stronę dla funkcji prostej o rozłącznych
nośnikach $g=\sum_jx_j\chi_{B_j}$ podział $A$ przez
$B_j$ daje równość. Dla dowolnego $g$ i prostej $s$
wykorzystujemy oszacowanie
$\big||\nu_g|(A)-|\nu_s|(A)\big|
\le|\nu_{g-s}|(A)\le\int_A\|g-s\|\,\mathrm d\mu$.
Zbieżność $s\to g$ w $L^1$ kończy dowód.
\end{proof}

\begin{lemat}[Różniczkowanie całki Bochnera]
\label{lem:bochner-differentiation-30}
Jeśli $g\in L^1((a,b),\lambda_1;X)$, to
\[
 \lim_{h\to0}\left\|
 \frac1h\int_{(t,t+h)}g(s)\,\mathrm d\lambda_1(s)-g(t)
 \right\|=0
\]
dla prawie każdego $t\in(a,b)$; dla $h<0$ całka
po przedziale jest rozumiana ze znakiem, tzn.
$\int_{(t,t+h)}=-\int_{(t+h,t)}$.
\end{lemat}
\begin{proof}
Wystarczy wykazać, że prawie każdy punkt jest
punktem Lebesgue'a funkcji $X$-wartościowej, czyli
średnia z $\|g(s)-g(t)\|$ na małych symetrycznych
przedziałach dąży do zera. Funkcje ciągłe
$\phi:[a,b]\to X$ są gęste w $L^1((a,b);X)$:
najpierw aproksymujemy $g$ funkcjami prostymi,
a funkcje charakterystyczne zbiorów mierzalnych
aproksymujemy w $L^1$ funkcjami ciągłymi,
korzystając z regularności miary Lebesgue'a.
Niech $\phi_j$ zbiega do $g$ w $L^1$ tak szybko, aby
$\sum_j2^j\|g-\phi_j\|_{L^1}<\infty$.
Dla $a_j(t)=\|g(t)-\phi_j(t)\|$ mamy
$\sum_j2^j a_j(t)<\infty$ dla prawie wszystkich $t$.
Słaba nierówność maksymalna z Rozdziału
\ref{ch:differentiation-25} daje dla ustalonego
$\varepsilon>0$ oszacowanie
\[
 \lambda_1\!\left(\left\{t:
 \sup_{r>0}\frac1{2r}\int_{(t-r,t+r)}a_j(s)
 \,\mathrm d\lambda_1(s)>\varepsilon\right\}\right)
 \le C\varepsilon^{-1}\|g-\phi_j\|_{L^1}.
\]
Przez $j\to\infty$ odpowiednie zbiory wyjątkowe
mają miarę dążącą do zera. Jednostajna ciągłość
$\phi_j$ i nierówność trójkąta pokazują więc, że
$\limsup_{r\downarrow0}\frac1{2r}\int_{(t-r,t+r)}
\|g(s)-g(t)\|\,\mathrm d\lambda_1(s)\le2\varepsilon$
poza zbiorem o dowolnie małej mierze.
Przy $\varepsilon\downarrow0$ otrzymujemy tezę.
Średnią po przedziale jednostronnym ograniczamy
przez dwukrotność średniej symetrycznej.
\end{proof}

\section{Własność Radona--Nikodyma przestrzeni Banacha}
\label{sec:rnp-30}
\begin{definicja}[Własność Radona--Nikodyma]
\label{def:RNP-30}
Przestrzeń Banacha $X$ ma \emph{własność Radona--Nikodyma}
(w skrócie RNP), jeśli dla każdej przestrzeni z miarą
skończoną $(\Omega,\mathcal A,\mu)$ oraz każdej
$X$-wartościowej miary wektorowej $\nu$ o skończonej
wariacji, spełniającej $\nu\ll\mu$, istnieje
$g\in L^1(\mu;X)$ takie, że
\begin{equation}\label{eq:RNP-30}
 \nu(A)=\int_Ag(\omega)\,\mathrm d\mu(\omega)
 \qquad(A\in\mathcal A).
\end{equation}
Funkcja $g$ jest wyznaczona jednoznacznie
$\mu$-prawie wszędzie.
\end{definicja}
\begin{proof}[Jednoznaczność]
Niech $g,h$ mają takie same całki po wszystkich
$A\in\mathcal A$. Funkcja $k=g-h$ ma istotnie
separowalny zbiór wartości. Istnieje przeliczalna
rodzina funkcjonałów $x_j^*\in X^*$ o normach co
najwyżej $1$, która normuje tę separowalną przestrzeń:
$\|x\|=\sup_j|x_j^*(x)|$ dla jej elementów
(konstruujemy ją przez wybranie funkcjonałów
normujących punkty przeliczalnego zbioru gęstego,
a następnie używamy ciągłości).
Dla każdego $j$ skalarny Radon--Nikodym i
$\int_A x_j^*(k)\,\mathrm d\mu=0$ dla każdego $A$
dają $x_j^*(k)=0$ prawie wszędzie. Po usunięciu
przeliczalnej sumy zbiorów zerowych otrzymujemy
$\|k\|=0$. Stąd $g=h$ prawie wszędzie.
\end{proof}

\begin{przyklad}[Przestrzenie skończenie wymiarowe]
Każda przestrzeń Banacha $X$ o skończonym wymiarze
ma RNP. Wybierając bazę $e_1,\ldots,e_d$,
zapisujemy $\nu=\sum_j\nu_je_j$.
Twierdzenie Radona--Nikodyma dla miar skalarnych
daje $\nu_j=g_j\mu$. Funkcja
$g=\sum_jg_je_j$ jest silnie mierzalna i całkowalna,
gdyż $\|g\|\le C\sum_j|g_j|$ dla pewnej stałej $C$.
Wtedy $\nu=g\mu$. Ta argumentacja nie działa
bezpośrednio przy nieskończenie wielu współrzędnych.
\end{przyklad}

\begin{lemat}[Przybliżanie miary przez skończone podziały]
\label{lem:rnp-partition-30}
Niech $(\Omega,\mathcal A,\mu)$ będzie przestrzenią
o mierze skończonej, a $\nu:\mathcal A\to X$ miarą
wektorową spełniającą $\|\nu(A)\|\le\mu(A)$
dla każdego $A\in\mathcal A$.
Dla każdego skończonego podziału mierzalnego
$\pi=\{A_1,\ldots,A_N\}$ definiujemy
\[
 g_\pi=\sum_{j:\mu(A_j)>0}
 \frac{\nu(A_j)}{\mu(A_j)}\chi_{A_j}
 \quad\in L^1(\mu;X).
\]
Wówczas $\|g_\pi\|\le1$ prawie wszędzie,
a następujące warunki są równoważne:
\begin{enumerate}[label=(\roman*)]
 \item istnieje $g\in L^1(\mu;X)$ z
       $\nu(A)=\int_Ag\,\mathrm d\mu$ dla każdego $A$;
 \item sieć $(g_\pi)$, uporządkowana przez
       rozdrobnienie skończonych podziałów,
       jest siecią Cauchy'ego w $L^1(\mu;X)$.
\end{enumerate}
Jeśli $\rho$ rozdrabnia $\pi$, to
$g_\pi=\mathbb E_\pi g_\rho$, gdzie
$\mathbb E_\pi$ oznacza uśrednienie na komórkach
$\pi$. Operator $\mathbb E_\pi$ jest kontrakcją
na $L^1(\mu;X)$.
\end{lemat}
\begin{proof}
Na komórce $A_j$ uśrednienie funkcji $h$ jest równe
$\mu(A_j)^{-1}\int_{A_j}h\,\mathrm d\mu$
(na komórkach zerowych przyjmujemy $0$).
Z nierówności trójkąta dla całki Bochnera mamy
$\|\mathbb E_\pi h\|_{L^1}\le\|h\|_{L^1}$.
Jeżeli $\rho$ rozdrabnia $\pi$, to
$\int_{A_j}g_\rho\,\mathrm d\mu=\nu(A_j)$,
co daje równanie z tezy.

Załóżmy (i). Wtedy $g_\pi=\mathbb E_\pi g$.
Dla $\varepsilon>0$ wybierzmy funkcję prostą
$s$ z $\|g-s\|_{L^1}<\varepsilon$ i podział
$\pi_0$, na którego komórkach $s$ jest stała.
Dla każdego $\pi$ rozdrabniającego $\pi_0$
zachodzi $\mathbb E_\pi s=s$, skąd
\[
 \|g_\pi-g\|_{L^1}
 \le\|\mathbb E_\pi(g-s)\|_{L^1}+\|s-g\|_{L^1}
 <2\varepsilon.
\]
Zatem sieć jest zbieżna, a więc Cauchy'ego.

Jeśli zachodzi (ii), zupełność $L^1(\mu;X)$
daje granicę $g$. Ustalmy $A\in\mathcal A$.
Dla każdego podziału $\pi$ rozdrabniającego
$\{A,\Omega\setminus A\}$ mamy
$\int_Ag_\pi\,\mathrm d\mu=\nu(A)$.
Przechodząc do granicy w $L^1$ otrzymujemy
$\int_Ag\,\mathrm d\mu=\nu(A)$.
\end{proof}

\begin{lemat}[Całka po krzywej z pochodną Banachowską]
\label{lem:FTC-Banach-30}
Niech $F:[a,b]\to X$ będzie Lipschitzowska i
różniczkowalna w normie $X$ dla prawie wszystkich
$t\in(a,b)$. Wtedy jej pochodna $F'$ ma silnie
mierzalnego przedstawiciela, jest istotnie ograniczona,
a dla wszystkich $a\le s<t\le b$ zachodzi
\[
 F(t)-F(s)=\int_{(s,t)}F'(r)\,\mathrm d\lambda_1(r).
\]
\end{lemat}
\begin{proof}
W punktach różniczkowalności $F'$ jest granicą
ilorazów różnicowych z dodatnimi wymiernymi
przyrostami. Każdy iloraz jest funkcją ciągłą na
odpowiednim przedziale, o wartościach w separowalnej
podprzestrzeni rozpiętej przez $F([a,b])$.
Przedłużamy $F'$ przez zero na zbiorze wyjątkowym;
punktowa granica funkcji silnie mierzalnych jest
silnie mierzalna. Ponadto $\|F'\|\le\operatorname{Lip}(F)$.
Niech $x^*\in X^*$. Funkcja skalarna
$x^*\circ F$ jest Lipschitzowska, a prawie wszędzie
$(x^*\circ F)'=x^*(F')$. Jednowymiarowe
Twierdzenie~\ref{lem:rademacher-line-28} daje
\[
 x^*(F(t)-F(s))
 =\int_{(s,t)}x^*(F'(r))\,\mathrm d\lambda_1(r)
 =x^*\!\left(\int_{(s,t)}F'(r)\,\mathrm d\lambda_1(r)\right).
\]
Ponieważ zachodzi to dla każdego $x^*$,
Hahn--Banach dowodzi równości wektorów.
\end{proof}

\begin{lemat}[Dzielenie miary bezatomowej]
\label{lem:nonatomic-divisibility-30}
Niech $\mu$ będzie skończoną miarą bezatomową.
Dla każdego mierzalnego $A$ i każdej liczby
$0\le t\le\mu(A)$ istnieje mierzalny $B\subseteq A$
taki, że $\mu(B)=t$. W konsekwencji dla $\eta>0$
każdy mierzalny $A$ można podzielić na skończenie
wiele zbiorów miary co najwyżej $\eta$.
\end{lemat}
\begin{proof}
Najpierw zauważmy, że każdy zbiór dodatniej miary
zawiera podzbiór o dowolnie małej dodatniej mierze.
Istotnie, jeśli $E$ ma dodatnią miarę, z braku atomów
istnieje $D\subset E$ z $0<\mu(D)<\mu(E)$.
Mniejszy z $D$ i $E\setminus D$ ma dodatnią miarę
nie większą niż $\mu(E)/2$. Iterując otrzymujemy
zbiory dodatniej miary mniejszej niż zadane $\eta$.

Dla $0<t<\mu(A)$ rozważmy częściowo uporządkowaną
przez inkluzję modulo zbiory zerowe rodzinę
$\mathcal C=\{B\subseteq A:B\in\mathcal A,\ \mu(B)\le t\}$.
Każdy łańcuch ma kres górny: jeśli $s$ jest kresem
górnym miar zbiorów łańcucha, wybieramy rosnący
(w porządku modulo zbiory zerowe) ciąg $B_j$ o
$\mu(B_j)\uparrow s$, a jego suma $B=\bigcup_jB_j$
stanowi ograniczenie górne całego łańcucha modulo
zbiory zerowe. Rzeczywiście, każdy zbiór łańcucha
o mierze mniejszej od $s$ mieści się modulo zbiory
zerowe w pewnym $B_j$, a każdy zbiór o mierze $s$
jest modulo zbiory zerowe równy $B$.
Z lematu Zorna istnieje element maksymalny $B$.
Jeśli $\mu(B)<t$, to $\mu(A\setminus B)>0$;
wybierając w $A\setminus B$ zbiór $C$ o mierze
$0<\mu(C)<t-\mu(B)$, otrzymalibyśmy większy
element $B\cup C\in\mathcal C$, sprzeczność.
Stąd $\mu(B)=t$. Przypadki $t=0$ i $t=\mu(A)$
są bezpośrednie. Dzieląc kolejno $A$ na części
miary $\eta$ aż do uzyskania reszty o mierze
co najwyżej $\eta$, dostajemy skończony podział.
\end{proof}

\begin{twierdzenie}[RNP a różniczkowalność krzywych Lipschitzowskich]
\label{thm:rnp-lipschitz-equiv-30}
Dla rzeczywistej lub zespolonej przestrzeni Banacha $X$
następujące warunki są równoważne:
\begin{enumerate}[label=(\roman*)]
 \item $X$ ma własność Radona--Nikodyma z
       Definicji~\ref{def:RNP-30};
 \item każda funkcja Lipschitzowska
       $F:[0,1]\to X$ jest różniczkowalna w normie $X$
       dla $\lambda_1$-prawie każdego $t\in(0,1)$.
\end{enumerate}
Wtedy każda taka $F$ spełnia wzór
$F(t)=F(0)+\int_{(0,t)}F'(s)\,\mathrm d\lambda_1(s)$.
\end{twierdzenie}
\begin{proof}
\emph{Kierunek (i)$\Rightarrow$(ii).}
Niech $F:[0,1]\to X$ ma stałą Lipschitza $L$.
Na algebrze skończonych sum rozłącznych przedziałów
półotwartych definiujemy
$\nu((s,t])=F(t)-F(s)$. Definicja jest skończenie
addytywna, ponieważ przy dzieleniu przedziału przyrosty
się teleskopują. Z nierówności Lipschitza wynika
$\|\nu(A)\|\le L\lambda_1(A)$ dla takich zbiorów.
Każdy zbiór mierzalny $A\subseteq[0,1]$ można
przybliżyć w mierze Lebesgue'a skończonymi sumami
przedziałów półotwartych. Przybliżające wartości
$\nu(A_j)$ są ciągiem Cauchy'ego w $X$:
$\|\nu(A_j)-\nu(A_k)\|\le L\lambda_1(A_j\triangle A_k)$.
Ich granica nie zależy od aproksymacji i określa
rozszerzenie $\nu$ na $\sigma$-algebrę Lebesgue'a z
$\|\nu(A)\|\le L\lambda_1(A)$. Dla rozłącznych
$A_j$ to oszacowanie kontroluje normę
$\nu(\bigcup_{j>N}A_j)$, zatem rozszerzenie jest
przeliczalnie addytywne. Ma skończoną wariację:
$|\nu|(A)\le L\lambda_1(A)$.

Na mocy RNP istnieje $g\in L^1(\lambda_1;X)$ takie,
że $\nu(A)=\int_Ag\,\mathrm d\lambda_1$.
Dla $t\in[0,1]$ otrzymujemy
$F(t)=F(0)+\int_{(0,t)}g\,\mathrm d\lambda_1$.
Lemat~\ref{lem:bochner-differentiation-30} daje
$F'=g$ prawie wszędzie.

\emph{Kierunek (ii)$\Rightarrow$(i): redukcja do części bezatomowej.}
Niech $\nu:\mathcal A\to X$ ma skończoną wariację
oraz $\nu\ll\mu_0$ dla skończonej miary $\mu_0$.
Wybierzmy miarę $\mu=|\nu|$. Jest skończona i
spełnia $\|\nu(A)\|\le\mu(A)$.
Jeśli $\nu=g\mu$ dla pewnego $g\in L^1(\mu;X)$,
to na mocy Lematu~\ref{lem:vector-variation-ac-30}
mamy $\mu\ll\mu_0$. Skalarne twierdzenie
Radona--Nikodyma daje $w=\mathrm d\mu/\mathrm d\mu_0$;
wtedy $gw\in L^1(\mu_0;X)$ reprezentuje $\nu$
względem $\mu_0$. Wystarczy zatem zbudować
gęstość względem $\mu=|\nu|$.

Atomy dodatniej miary skończonej $\mu$ można
wybrać parami rozłącznie, a jest ich co najwyżej
przeliczalnie wiele, licząc z dokładnością do zbiorów
zerowych. Istotnie, dwa atomy o dodatnim przecięciu
są równe modulo zbiór zerowy. Dla każdego $k$
istnieje co najwyżej skończenie wiele rozłącznych
atomów miary większej niż $1/k$; ich sumy, po
przejściu po wszystkich $k$, obejmują część atomową.
Po odjęciu tej części pozostaje miara bezatomowa.
Na atomie $A$ funkcja
$\nu$ przyjmuje wartości proporcjonalne do $\mu$,
ponieważ każda mierzalna część atomu ma miarę $0$
lub $\mu(A)$; gęstość wynosi tam
$\nu(A)/\mu(A)$. Suma takich części daje funkcję
$g_{\mathrm{at}}\in L^1(\mu;X)$, gdyż
$\sum_A\|\nu(A)\|\le\sum_A\mu(A)\le\mu(\Omega)$.
Pozostaje część bezatomowa. Zakładamy więc dalej,
że $\mu$ jest bezatomowa i $\mu(\Omega)=M>0$.

\emph{Przypuszczenie o braku gęstości i wybór podziałów.}
Załóżmy przeciwnie, że $\nu$ nie ma gęstości
względem $\mu$. Z Lematu~\ref{lem:rnp-partition-30}
wynika, że sieć funkcji $g_\pi$ nie jest Cauchy'ego.
Istnieje $\varepsilon>0$ takie, że dla każdego
podziału $\pi_0$ znajdują się jego rozdrobnienia
$\alpha,\beta$ z
$\|g_\alpha-g_\beta\|_{L^1}\ge\varepsilon$.
Niech $\rho$ będzie wspólnym rozdrobnieniem
$\alpha$ i $\beta$. Ponieważ uśrednienia są
kontrakcjami i $\mathbb E_\alpha g_\rho=g_\alpha$,
$\mathbb E_\beta g_\rho=g_\beta$,
$\mathbb E_\alpha g_{\pi_0}=g_{\pi_0}$ oraz
$\mathbb E_\beta g_{\pi_0}=g_{\pi_0}$,
otrzymujemy
$\varepsilon\le2\|g_\rho-g_{\pi_0}\|_{L^1}$.
Na mocy Lematu~\ref{lem:nonatomic-divisibility-30}
każdy mierzalny zbiór można rozciąć na skończenie
wiele części o arbitralnie małej mierze. Jeśli dalej rozdrobnimy
$\rho$ do $\pi$, to ponownie przez kontrakcję
$\|g_\pi-g_{\pi_0}\|_{L^1}
\ge\|g_\rho-g_{\pi_0}\|_{L^1}$.
Możemy więc zbudować zagnieżdżone podziały
$\pi_1\prec\pi_2\prec\cdots$ takie, że
\begin{equation}\label{eq:rnp-partition-jumps-30}
 \max_{A\in\pi_n}\mu(A)\longrightarrow0,
 \qquad
 \|g_{\pi_{n+1}}-g_{\pi_n}\|_{L^1}
 \ge\varepsilon/2\quad(n\ge1).
\end{equation}
Zbiory o mierze zero pomijamy.

\emph{Kodowanie podziałów jako przedziałów.}
Każdej komórce $A\in\pi_1$ przypisujemy
półotwarty przedział $I_A\subseteq[0,M]$ o
długości $\mu(A)$, tak aby przedziały te
w kolejności wypełniały $[0,M]$.
Jeśli komórka $A$ jest podzielona na
$A_1,\ldots,A_r$ przez $\pi_2$, dzielimy
$I_A$ na kolejne podprzedziały długości
$\mu(A_1),\ldots,\mu(A_r)$. Powtarzamy to
rekurencyjnie. W ten sposób dostajemy zagnieżdżone
podziały przedziału, których maksymalna długość
komórki zmierza do zera. Niech $h_n:[0,M]\to X$
będzie funkcją stałą na $I_A$ dla $A\in\pi_n$
o wartości $\nu(A)/\mu(A)$. Z konstrukcji
\[
 \|h_n\|\le1\ \text{prawie wszędzie},\qquad
 \int_{I_A}h_n(t)\,\mathrm d\lambda_1(t)=\nu(A).
\]
Ponadto
\[
 \|h_{n+1}-h_n\|_{L^1([0,M];X)}
 =\|g_{\pi_{n+1}}-g_{\pi_n}\|_{L^1(\mu;X)}.
\]
Ostatnia równość wynika z porównania stałych
wartości na odpowiednich komórkach wspólnych
rozdrobnień, mających te same miary.

Niech $D$ będzie zbiorem końców wszystkich
przedziałów z wszystkich poziomów. Jest on
gęsty w $[0,M]$. Na $D$ definiujemy $F(0)=0$
i określamy $F(t)$ jako sumę $\nu(A)$ po
komórkach leżących na lewo od $t$, na dowolnym
poziomie, na którym $t$ jest końcem przedziału.
Definicja jest spójna dzięki skończonej addytywności
$\nu$; dla $s<t$ z $D$ mamy
$\|F(t)-F(s)\|\le\mu(\text{suma odpowiednich komórek})
=t-s$. Zatem $F$ przedłuża się jednoznacznie
do $1$-Lipschitzowskiej funkcji na $[0,M]$.
Dla każdej komórki $I_A=[c,d]$ zachodzi
\begin{equation}\label{eq:rnp-interval-average-30}
 F(d)-F(c)=\nu(A)
 =\int_{I_A}h_n(t)\,\mathrm d\lambda_1(t).
\end{equation}

\emph{Różniczkowanie wymusza sprzeczność.}
Warunek (ii), po liniowym przeskalowaniu $[0,M]$
do $[0,1]$, zapewnia istnienie $F'$ prawie wszędzie.
Lemat~\ref{lem:FTC-Banach-30} daje
$F(d)-F(c)=\int_{(c,d)}F'(t)\,\mathrm d\lambda_1(t)$.
Porównanie z \eqref{eq:rnp-interval-average-30}
pokazuje, że na każdej komórce $I_A$ funkcja $h_n$
jest średnią $F'$; inaczej mówiąc,
$h_n=\mathbb E_{\mathcal I_n}F'$,
gdzie $\mathcal I_n$ oznacza skończony podział
$[0,M]$ na opisane przedziały.
Ponieważ ich maksymalna długość dąży do zera,
$h_n\to F'$ w $L^1([0,M];X)$. Istotnie, dla ciągłej
$\phi:[0,M]\to X$ średnie po przedziałach różnią się
od $\phi$ o nie więcej niż jej moduł ciągłości
przy maksymalnej długości komórki. Dla ogólnej
$F'\in L^1$ przybliżamy ją funkcją ciągłą i
korzystamy z kontrakcji operatorów uśredniających.
Otrzymana zbieżność przeczy nierównościom
\eqref{eq:rnp-partition-jumps-30}. Zatem $\nu$
ma gęstość względem $|\nu|$, a stąd, jak wykazano
na początku, względem $\mu_0$.
\end{proof}

\section{Przykłady: przestrzenie z RNP i bez RNP}
\label{sec:rnp-examples-30}
\begin{twierdzenie}[Przestrzenie Hilberta mają RNP]
\label{thm:hilbert-rnp-30}
Każda rzeczywista lub zespolona przestrzeń Hilberta $H$
ma własność Radona--Nikodyma.
\end{twierdzenie}
\begin{proof}
Niech $\nu:\mathcal A\to H$ ma skończoną wariację
i niech $\mu=|\nu|$. Dla podziału $\pi$ funkcja
$g_\pi$ z Lematu~\ref{lem:rnp-partition-30} spełnia
$\|g_\pi\|_{L^\infty}\le1$, zatem jest ograniczona
w przestrzeni Hilberta $L^2(\mu;H)$:
$\|g_\pi\|_{L^2}^2\le\mu(\Omega)$.
W domkniętej kuli tej przestrzeni istnieje
słabo zbieżna podsieć $(g_{\pi_\alpha})$,
gdyż kula przestrzeni Hilberta jest słabo zwarta.
Niech $g\in L^2(\mu;H)$ będzie jej granicą.
Dla każdego $A\in\mathcal A$ i $h\in H$ funkcja
$\chi_A h$ należy do $L^2(\mu;H)$. Jeśli
$\pi_\alpha$ rozdrabnia $\{A,\Omega\setminus A\}$,
mamy $\int_Ag_{\pi_\alpha}\,\mathrm d\mu=\nu(A)$.
Zatem słaba zbieżność w $L^2$ daje
\[
 \left\langle\int_Ag\,\mathrm d\mu,h\right\rangle_H
 =\lim_\alpha\left\langle
 \int_Ag_{\pi_\alpha}\,\mathrm d\mu,h\right\rangle_H
 =\langle\nu(A),h\rangle_H.
\]
Dowolność $h$ dowodzi $\nu(A)=\int_Ag\,\mathrm d\mu$.
Ponieważ $\mu$ jest skończona, $L^2(\mu;H)$
zanurza się w $L^1(\mu;H)$.
Dla pierwotnej miary dominującej $\mu_0$
przechodzimy z $\mu=|\nu|$ do $\mu_0$ dokładnie jak
w dowodzie Twierdzenia~\ref{thm:rnp-lipschitz-equiv-30}.
\end{proof}

\begin{przyklad}[Dlaczego przestrzeń $L^1$ nie ma RNP]
\label{ex:L1-no-RNP-30}
Niech $X=L^1([0,1],\lambda_1)$ oraz
$F:[0,1]\to X$ będzie określona przez
$F(t)=\chi_{[0,t]}$ (jako klasa funkcji w zmiennej $s$).
Dla $t>u$ mamy $\|F(t)-F(u)\|_{L^1}=t-u$,
a więc $F$ jest $1$-Lipschitzowska.
W żadnym punkcie $t\in(0,1)$ nie istnieje jednak
pochodna w normie $L^1$. Istotnie, dla $h>0$
iloraz różnicowy to
$h^{-1}\chi_{(t,t+h]}$, którego norma wynosi $1$.
Gdyby istniała granica $v\in L^1$, to zbieżność
w $L^1$ pewnego podciągu dawałaby zbieżność
prawie wszędzie do $v$. Dla każdego $s\ne t$
ilorazy różnicowe są jednak ostatecznie równe $0$.
Zatem $v=0$ prawie wszędzie, sprzecznie z
$\|h^{-1}\chi_{(t,t+h]}\|_{L^1}=1$.
Z Twierdzenia~\ref{thm:rnp-lipschitz-equiv-30}
wynika, że $L^1([0,1])$ nie ma RNP.

Równoważnie, miara wektorowa
$\nu(A)=\chi_A\in L^1([0,1])$ spełnia
$\|\nu(A)\|_{L^1}=\lambda_1(A)$ i
$|\nu|=\lambda_1$, ale \emph{nie} ma gęstości
Bochnera $g$ względem $\lambda_1$.
Gdyby miała, jej funkcja pierwotna $F(t)=\nu([0,t])$
musiałaby być różniczkowalna prawie wszędzie.
\end{przyklad}

\begin{przyklad}[Przestrzeń funkcji ciągłych bez RNP]
Niech $X=C([0,1])$ z normą supremum i
$F(t)(s)=|s-t|$ dla $s,t\in[0,1]$.
Wtedy $\|F(t)-F(u)\|_\infty\le|t-u|$.
Dla ustalonego $t\in(0,1)$ skalarne ilorazy różnicowe
przy $h\downarrow0$ dążą punktowo (dla $s\ne t$)
do funkcji równej $+1$ dla $s<t$ oraz $-1$ dla $s>t$.
Ta funkcja ma skok, więc nie może być jednostajną
granicą funkcji ciągłych na $[0,1]$.
Nie istnieje zatem pochodna $F'(t)$ w normie supremum,
a $C([0,1])$ nie ma RNP.
\end{przyklad}

\begin{uwaga}[Nie należy mylić dwóch twierdzeń Rademachera]
Twierdzenie~\ref{thm:rademacher-28} dowodzi
różniczkowalności odwzorowań Lipschitzowskich
$U\subseteq\mathbb R^n\to\mathbb R^m$ prawie
wszędzie, bez dodatkowych założeń. Dla ogólnej
przestrzeni Banacha $X$ nawet odwzorowanie
$[0,1]\to X$ może być nigdzie nieróżniczkowalne.
Twierdzenie~\ref{thm:rnp-lipschitz-equiv-30}
wykazuje, że dokładną przeszkodą jest brak RNP.
W szczególności nie wolno wnioskować o RNP
z samego istnienia skalarnej pochodnej
Radona--Nikodyma dla $x^*\nu$ przy każdym
$x^*\in X^*$: te skalarne gęstości nie muszą
składać się w silnie mierzalną gęstość $X$-wartościową.
\end{uwaga}

\section{Miary wektorowe w analizie i geometrii}
\label{sec:vector-applications-30}
\begin{przyklad}[Miara gradientowa funkcji $BV$]
Jeśli $u\in BV(\Omega)$, to z
Twierdzenia~\ref{thm:BV-measure-29} dystrybucyjna
pochodna $Du=(D_1u,\ldots,D_nu)$ jest skończoną
miarą wektorową o wartościach w $\mathbb R^n$.
Dla $u\in BV_{\mathrm{loc}}(\Omega)$ stosujemy
poniższe wnioski o gęstościach lokalnie, po ograniczeniu
do części o skończonej wariacji.
Z RNP tej przestrzeni i twierdzenia
Radona--Nikodyma otrzymujemy gęstość
$\sigma_u=\mathrm dDu/\mathrm d|Du|$ o normie $1$
$|Du|$-prawie wszędzie. Dla funkcji gładkiej
$Du=\nabla u\,\lambda_n$, natomiast dla funkcji
charakterystycznej zbioru o skończonym obwodzie
miara $D\chi_E$ może być osobliwa względem
$\lambda_n$ i skoncentrowana na brzegu miarowym.
To wyjaśnia, dlaczego naturalną miarą dominującą
jest $|Du|$, a nie zawsze $\lambda_n$.
\end{przyklad}

\begin{wniosek}[Przestrzenie Banacha i całki z krzywych]\label{idx:30:wniosek:15}

Jeśli $X$ ma RNP i $F:[a,b]\to X$ jest
Lipschitzowska, to istnieje $g\in L^\infty((a,b);X)$
takie, że $F(t)-F(a)=\int_{(a,t)}g\,\mathrm d\lambda_1$
i $F'=g$ prawie wszędzie.
\end{wniosek}
\begin{proof}
Po przeskalowaniu przedziału stosujemy
Twierdzenie~\ref{thm:rnp-lipschitz-equiv-30} i
Lemat~\ref{lem:FTC-Banach-30}.
\end{proof}

\section{Granice zakresu i dalsze kierunki}
Teoria miar wektorowych ma również zagadnienia,
których nie potrzebowaliśmy do dowodu RNP:
całkę Pettisa, miary o skończonej półwariacji,
operatory słabo zwarte oraz twierdzenie
Bartle'a--Dunforda--Schwartza.
Są one naturalną kontynuacją w analizie funkcjonalnej.
W tym rozdziale udowodniliśmy od podstaw
charakteryzację RNP przez różniczkowalność krzywych
Lipschitzowskich. Nie zakładaliśmy ani twierdzenia
o dentowalności zbiorów wypukłych, ani ogólnego
twierdzenia o zbieżności martyngałów Banachowskich.
